\documentclass[conference]{IEEEtran}
\usepackage[T1]{fontenc}
\usepackage{cite}
\usepackage{graphicx}
\usepackage{url}
\usepackage{booktabs}
\usepackage{tikz}
\usetikzlibrary{arrows.meta,positioning,fit,calc}
\usepackage{xcolor}
\usepackage{amsmath}

\definecolor{scholarblue}{RGB}{35,82,124}
\definecolor{scholargray}{RGB}{242,245,248}
\tikzset{
  block/.style={draw=scholarblue, rounded corners=2pt, fill=scholargray,
    align=center, minimum width=1.55cm, minimum height=0.58cm, font=\scriptsize},
  arrow/.style={-Latex, draw=scholarblue, line width=0.55pt}
}

\title{ScholarLM: An Interactive PDF Study Workspace with\\Retrieval-Augmented Assistance and Canvas-Based Knowledge Mapping}

\author{\IEEEauthorblockN{T Aditya Ananda Sai}
\IEEEauthorblockA{Term-3 Project Report\\
Indian Institute of Technology Guwahati\\
taa.sai@op.iitg.ac.in}}

\begin{document}
\maketitle

\begin{abstract}
ScholarLM is an interactive PDF study workspace that combines document reading, an infinite annotation canvas, retrieval-augmented question answering, visual explanations, and a knowledge graph in one browser application. The system is designed to reduce the context switching that occurs when learners move between a PDF viewer, a notes application, and an AI assistant. A Bun and Hono backend extracts page-aware PDF text, tables, and selected visual context; chunks and embeds the content; and stores documents, canvas state, notes, embeddings, and graph data in SQLite. A React frontend uses PDF.js and tldraw to place the document and user reasoning on a persistent canvas. OpenRouter provides configurable language, vision, and embedding models, while Kokoro or browser speech provides read-aloud support. This paper presents the architecture, grounding strategy, deterministic equation-plotting pipeline, persistence model, and functional validation of the current implementation.
\end{abstract}

\begin{IEEEkeywords}
PDF learning, retrieval-augmented generation, interactive canvas, semantic search, knowledge graph, handwriting recognition, educational software
\end{IEEEkeywords}

\section{Introduction}
Academic reading commonly requires several tools at once: a document reader for source material, a notes application for annotation, a search interface for locating concepts, and an AI assistant for explanation. Moving information between these tools interrupts reading and makes it difficult to preserve the relationship between a note and the exact page or passage that motivated it.

Large language models can make technical material easier to explore, but a generic conversational interface is not sufficient for document study. Answers should be grounded in the open document, citations should lead back to source pages, and user-created explanations should remain attached to the learner's workspace. Visual material introduces an additional challenge because important evidence may occur in diagrams, tables, equations, or handwriting rather than in the PDF text layer.

ScholarLM addresses these requirements through a document-centered workspace. A PDF is displayed inside a persistent tldraw canvas, where the learner can draw, type, highlight, create sticky notes, and place generated explanations. The same workspace exposes Explain, Ask, Search, Notes, and Graph panels. Its main contributions are:

\begin{itemize}
\item an integrated PDF and infinite-canvas study environment;
\item page-aware ingestion with chunking, embeddings, table context, repeated-header filtering, and selected visual context;
\item grounded document question answering with source citations and refusal when evidence is insufficient;
\item a graph connecting documents, concepts, notes, explanations, handwriting, and canvas objects;
\item a restricted, deterministic plotting path for supported equations recognized from text or handwriting; and
\item persistence and deletion workflows that preserve or remove affiliated data across SQLite and browser recovery storage.
\end{itemize}

\section{Related Work}
Retrieval-augmented generation (RAG) combines information retrieval with language-model generation so that a response can use an external knowledge source rather than relying only on pretrained parameters \cite{rag}. Semantic vector search is particularly useful when a learner's wording differs from the wording used in the source document. PDF.js provides browser-side access to PDF rendering and text layers \cite{pdfjs}, while tldraw provides an extensible infinite-canvas interaction model \cite{tldraw}.

Knowledge graphs offer a complementary representation. A vector index answers similarity-oriented questions, whereas a graph makes entities and their relationships navigable. ScholarLM combines both: semantic retrieval supplies evidence for Ask PDF and Search, while graph nodes connect source documents to concepts and saved user artifacts. Unlike a conventional chat interface, the generated response is treated as a workspace object that can be revisited, moved, and linked to its source.

\section{System Architecture}
The overall data flow is shown in Fig.~\ref{fig:architecture}. The frontend owns document viewing and canvas interaction. The backend provides API routes, persistence, ingestion, retrieval, model calls, speech generation, and graph construction.

\begin{figure}[t]
\centering
\includegraphics[width=\columnwidth]{architecture_diagram.jpg}
\caption{Architecture of the ScholarLM system.}
\label{fig:architecture}
\end{figure}
\subsection{Document Ingestion}
When a PDF is uploaded, it is stored in the backend upload directory and recorded in SQLite. The ingestion pipeline progresses through extracting, chunking, embedding, graphing, and ready states. PDF parsing records page numbers, text, tables, and candidate visual pages. Pages with images, sparse text, or indicators such as ``chart'' and ``diagram'' can be rendered for additional visual context. Repeated watermarks and headers are filtered before indexing.

The normalized pages are chunked with a maximum size of 1800 characters and an overlap of 120 characters. Chunk boundaries prefer paragraphs, sentences, and word boundaries. Each chunk retains its source page and receives a document embedding. Embeddings are stored in SQLite and compared with normalized query vectors using cosine similarity. Sticky notes and saved explanations are indexed alongside PDF chunks, enabling a single search surface for both source material and learner-generated knowledge.

\subsection{Interactive Workspace}
The browser workspace renders the PDF and a tldraw canvas together. The user can select text from the PDF, lasso handwriting, paste or upload a screenshot, draw annotations, and create movable notes. Explanations can be saved back to the canvas as expandable sticky shapes. A standalone canvas can also be created before a PDF is attached. Canvas snapshots are saved in SQLite and backed up in browser storage so normal navigation and reloads do not discard work.

\section{Grounded Assistance}
\subsection{Ask PDF and Search}
For a document question, ScholarLM embeds the query and retrieves a small set of relevant page-aware chunks. The selected excerpts are supplied to the configured model with source identifiers such as S1 and S2. The generation prompt requires every factual claim to cite one or more supplied sources. If the retrieved evidence is insufficient, the system returns a refusal instead of filling the gap with outside knowledge. Cited sources are retained in the response so that selecting a citation can navigate back to the originating page.

The answer is cached using a content version derived from the document, chunks, and indexed stickies. Exact and sufficiently similar questions can therefore reuse a valid answer, while edits or re-indexing invalidate stale results. Search combines vector similarity with lexical matching and returns PDF passages, explanations, notes, handwriting, or concepts with navigation metadata.

\subsection{Explanation and Vision}
The Explain workflow supports selected PDF text, canvas text, handwriting, and screenshots. Text explanations are streamed when possible. Visual requests use the configured vision model and return a structured result containing an explanation, an optional recognized equation, and optional plot data. The prompt explicitly asks the model to state uncertainty rather than invent unreadable symbols. Explanations can be simplified or regenerated, and unchanged selections can use cached explanation records.

\section{Knowledge Graph and Equation Support}
After indexing, the backend asks the model to identify concepts and meaningful relationships from the document chunks. A graph response is validated before being stored: concepts have labels, descriptions, and page references, while edges refer to existing concept labels. The frontend visualizes document-scoped and library-wide graphs using graphology and Sigma. Manual groups and edges allow the learner to organize related documents or concepts. Graph nodes can point back to PDF pages, canvas shapes, sticky notes, explanations, and handwriting.

Equation support deliberately separates visual recognition from computation. The model may transcribe a selected equation and explain it, but plot points are generated by a restricted parser rather than by executing model-produced code. The current implementation supports explicit functions such as $y=\sin(x)$ and $y=x^2$, powers, roots, common trigonometric functions, discontinuities such as $y=1/x$, and origin-centered circles such as $x^2+y^2=1$. Unsupported expressions are rejected. Generated graphs remain linked to their source shapes and are invalidated when the source equation changes.

\section{Implementation}
The backend is a Bun and Hono TypeScript service using Bun's SQLite interface, pdf-parse, pdf-lib, Sharp, and Kokoro JS. The frontend is a React and TypeScript application using PDF.js, react-pdf, tldraw, graphology, Sigma, and Framer Motion. OpenRouter is used through configurable endpoints for generation, vision, and embeddings. The default configuration supports a free routed model and separately configurable vision and embedding models.

The persistence layer includes tables for documents, page text, chunks, embeddings, notes, explanations, graph concepts, graph edges, manual groups, speech-cache entries, and provider telemetry. Database startup performs schema initialization and lightweight migrations. Destructive operations remove affiliated files, indexes, graph edges, explanations, and browser recovery data so deleted study material does not reappear after reload.

The application exposes separate routes for documents, canvases, notes, explanations, retrieval, search, graph operations, and speech. Provider failures are retried for short transient errors when no output has been emitted. Kokoro speech is cached and streamed when available; browser speech is used as a fallback.

\section{Models and Fallback Strategy}
ScholarLM separates language generation, visual understanding, embeddings, and speech synthesis so that each capability can be configured independently. The model names below are the defaults declared in the project configuration; they can be changed through environment variables without modifying the application code.

\begin{table}[t]
\caption{Models and fallback behavior in ScholarLM}
\label{tab:models}
\centering
\scriptsize
\begin{tabular}{p{0.21\columnwidth} p{0.39\columnwidth} p{0.25\columnwidth}}
\toprule
\textbf{Capability} & \textbf{Primary configuration} & \textbf{Fallback}\\
\midrule
Text generation & \texttt{google/gemma-4-26b-a4b-it:free} & OpenRouter provider fallback or another configured routing model.\\
Vision and handwriting & \texttt{google/gemma-4-26b-a4b-it:free} & OpenRouter provider fallback when the selected provider is unavailable.\\
Document embeddings & \texttt{nvidia/llama-nemotron-embed-vl-1b-v2:free} & Re-indexing can be retried with the configured embedding provider; unavailable embeddings disable semantic retrieval until indexing succeeds.\\
Speech synthesis & Local Kokoro JS model & Browser SpeechSynthesis API when Kokoro cannot load or decode audio.\\
Transient provider errors & Current OpenRouter operation and selected model & Short retry with exponential backoff, provided that no response tokens have already been emitted.\\
\bottomrule
\end{tabular}
\end{table}

The current primary language model is Gemma 4 26B, while the embedding model is NVIDIA Llama Nemotron Embed VL 1B v2. Additional language-model routing candidates can be ordered through the \texttt{OPENROUTER\_ROUTING\_MODELS} configuration. OpenRouter provider fallback is enabled for embedding requests, and the backend records provider operations and failures for diagnosis. The application also preserves the active selection during failed explanation requests and reuses cached explanations or grounded answers when their source-content version is unchanged. These behaviors provide graceful degradation without silently substituting unsupported evidence.
\section{Functional Validation}
The repository documentation records a completed feature baseline and a manual clean-data workflow. Table~\ref{tab:validation} summarizes the implemented behaviors. These are functional checks rather than a quantitative benchmark; latency, retrieval accuracy, transcription accuracy, and grounded-answer rate remain evaluation tasks for a controlled test set.

\begin{table}[t]
\caption{Functional validation of the current ScholarLM implementation}
\label{tab:validation}
\centering
\scriptsize
\begin{tabular}{p{0.29\columnwidth} p{0.58\columnwidth}}
\toprule
\textbf{Test area} & \textbf{Observed capability}\\
\midrule
PDF ingestion & Page-aware extraction, table handling, chunking, embeddings, and status tracking are implemented.\\
Grounded QA & Retrieved passages are cited; unsupported questions use an evidence-based refusal.\\
Canvas persistence & PDF-linked and standalone tldraw canvases survive normal navigation and reload.\\
Explanation workflow & Text, handwriting, and screenshot selections can be explained, cached, and saved as stickies.\\
Search & PDF chunks and saved stickies share a semantic search path with source navigation.\\
Knowledge graph & Concepts, relationships, manual groups, and source-linked graph navigation are implemented.\\
Math plotting & Supported equations are parsed deterministically; unsupported expressions are rejected.\\
Speech & Kokoro generation, caching, playback controls, and browser fallback are available.\\
Deletion & Documents, notes, graphs, indexes, files, and recovery copies have cascade cleanup paths.\\
\bottomrule
\end{tabular}
\end{table}

The principal engineering result is an integrated study loop: upload a document, read and annotate it, ask a cited question, explain a selected passage, save the result as a note, and follow relationships through the graph. The system also keeps the user's current selection and inspector state available across normal navigation, which is important for maintaining context during exploration.

\section{Limitations}
The current build is primarily a local demonstration and is not a hosted multi-user service. It requires a Bun-compatible runtime, writable local storage, and an OpenRouter API key for model-backed features. There is no authentication, per-user isolation, quota management, or cloud synchronization.

Retrieval quality depends on PDF text extraction, chunk boundaries, embedding quality, and the quality of the retrieved evidence. Image-only PDFs, complex multi-column layouts, tables, and equations may require OCR or improved layout analysis. Visual recognition and handwriting transcription depend on image quality and model behavior. The knowledge graph is model-generated and therefore requires additional validation for duplicate concepts, malformed relationships, and provenance. The current math parser intentionally covers a limited language and does not provide arbitrary symbolic algebra.

The project has more manual verification than automated test coverage. A repeatable evaluation set is still needed for top-page retrieval accuracy, grounded versus unsupported answers, graph quality, equation transcription, cold and warm latency, speech fallback, and deletion during active ingestion.

\section{Future Work}
Future work includes OCR fallback, exact passage highlighting, retrieval confidence displays, regression tests for ingestion and retrieval, improved table and multi-column handling, and controlled evaluation on fixed question sets. The canvas experience can be extended with an explicit Math mode, automatic grouping of nearby strokes, equation correction, multiple equations on one graph, domain controls, tracing, and a richer equation panel.

Further engineering work should add numbered database migrations, automatic backups, recovery windows for destructive actions, conflict handling for multiple browser tabs, accessibility verification, mobile-specific layout, reduced-motion support, and security controls for hosted deployment. Multi-user collaboration and cloud synchronization are intentionally deferred until the local workflow is stable and its evaluation evidence is stronger.

\section{Conclusion}
ScholarLM demonstrates how document rendering, an infinite annotation canvas, semantic retrieval, grounded language-model assistance, visual analysis, deterministic equation plotting, speech, and knowledge graphs can be combined into a single learning workspace. Its central design principle is to keep the source PDF and the learner's reasoning in the same persistent context. The current implementation provides a practical local foundation for reading and exploring technical PDFs, while its explicit limitations define a clear path toward stronger evaluation, accessibility, reliability, and hosted deployment.

\begin{thebibliography}{00}
\bibitem{rag} P. Lewis et al., ``Retrieval-Augmented Generation for Knowledge-Intensive NLP Tasks,'' in \emph{Advances in Neural Information Processing Systems}, 2020.
\bibitem{pdfjs} Mozilla Foundation, ``PDF.js: A general-purpose, web standards-based platform for parsing and rendering PDFs,'' [Online]. Available: \url{https://mozilla.github.io/pdf.js/}
\bibitem{tldraw} tldraw, ``tldraw: Infinite canvas SDK for React,'' [Online]. Available: \url{https://tldraw.dev/}
\bibitem{bun} Oven, ``Bun documentation,'' [Online]. Available: \url{https://bun.sh/docs}
\bibitem{hono} Hono, ``Hono web framework documentation,'' [Online]. Available: \url{https://hono.dev/}
\bibitem{openrouter} OpenRouter, ``OpenRouter API documentation,'' [Online]. Available: \url{https://openrouter.ai/docs}
\bibitem{sqlite} SQLite, ``SQLite documentation,'' [Online]. Available: \url{https://sqlite.org/docs.html}
\end{thebibliography}

\end{document}








